%============================ BASE DE DATOS (ENTREGA 2) =====================
\chapter{Diseño de la base de datos}
\label{cap:basedatos}

\section{Estrategia de persistencia}
Por decisión de diseño ---priorizando la portabilidad, la ausencia de dependencias nativas y la facilidad de despliegue en cualquier entorno, incluida la nube--- EILOR implementa su capa de datos mediante \emph{persistencia en archivos JSON} gestionada por el módulo \code{lib/store}. Esta aproximación, adecuada para el volumen y la naturaleza del prototipo, se comporta como una base de datos documental ligera: cada colección se materializa en un archivo, con escritura atómica mediante archivo temporal y renombrado. La Sección~\ref{sec:migracion} describe la ruta de migración a un motor relacional para escenarios de mayor escala.

\section{Modelo entidad-relación}
La Figura~\ref{fig:er} presenta el modelo conceptual. Las entidades \textbf{Dispositivo}, \textbf{MuestraTemporal} y \textbf{Registro} son independientes en el modelo documental actual, pero se relacionan lógicamente con la entidad \textbf{PuntoAcceso} (el SoftAP activo) y con la sesión de captura.

\begin{figure}[H]
\centering
\begin{tikzpicture}[
  ent/.style={rectangle,draw,very thick,rounded corners=2pt,align=center,
              minimum width=3.0cm,minimum height=1.0cm,fill=eilorcyan!8,font=\small},
  rel/.style={diamond,aspect=1.6,draw,thick,align=center,font=\scriptsize,fill=eiloremerald!10,inner sep=1pt},
  arr/.style={-{Latex[length=2mm]},thick},
  attr/.style={font=\scriptsize\itshape,align=center}
]
\node[ent] (dev) {Dispositivo};
\node[ent,right=3.2cm of dev] (ap) {PuntoAcceso (SoftAP)};
\node[ent,below=2.0cm of dev] (ts) {MuestraTemporal};
\node[ent,below=2.0cm of ap] (reg) {Registro (invitado)};

\node[rel] (r1) at ($(dev)!0.5!(ap)$) {detecta};
\node[rel] (r2) at ($(ap)!0.5!(reg)$) {genera};

\draw[arr] (dev) -- (r1); \draw[arr] (r1) -- (ap);
\draw[arr] (ap) -- (r2); \draw[arr] (r2) -- (reg);
\draw[arr,dashed] (dev) -- node[attr,left]{agrega} (ts);

\node[attr,above=0.1cm of dev] {mac (PK), rssi, channel, vendor, evilTwin};
\node[attr,above=0.1cm of ap] {apSsid, apChannel, apActive};
\node[attr,below=0.1cm of ts] {t (PK), wifi, ble, active, rate};
\node[attr,below=0.1cm of reg] {id (PK), nombre, apellido, telefono, ts};
\end{tikzpicture}
\caption{Modelo entidad-relación conceptual de EILOR.}
\label{fig:er}
\end{figure}

\section{Esquema físico (archivos JSON)}
El módulo \code{lib/store} define tres archivos bajo el directorio \code{data/}, descritos en la Tabla~\ref{tab:archivos}.

\begin{table}[H]
\centering
\caption{Archivos de persistencia}
\label{tab:archivos}
\footnotesize
\begin{tabularx}{\linewidth}{@{}l l X@{}}
\toprule
\textbf{Archivo} & \textbf{Estructura} & \textbf{Contenido / política} \\
\midrule
\code{state.json} & Objeto & Instantánea de dispositivos y total de paquetes; \emph{snapshot} cada 30~s \\
\code{timeseries.json} & Arreglo & Muestras por minuto; retención máx. 720 (\code{MAX\_TIMESERIES}) \\
\code{registrations.json} & Arreglo & Altas de invitados; retención máx. 5000 (\code{MAX\_REGISTRATIONS}) \\
\bottomrule
\end{tabularx}
\end{table}

El Listado~\ref{lst:store} muestra la escritura atómica y el modelo de retención de la serie temporal.

\begin{lstlisting}[style=js,caption={Escritura atómica y retención (lib/store.js)},label={lst:store}]
function saveJSON(file, data) {
  try {
    ensureDir();
    const tmp = file + '.tmp';                 // escritura atomica:
    fs.writeFileSync(tmp, JSON.stringify(data)); // 1) escribe temporal
    fs.renameSync(tmp, file);                   // 2) renombra (atomico)
    return true;
  } catch (e) { return false; }
}

function saveTimeseries(series) {
  return saveJSON(HISTORY_FILE, series.slice(-MAX_TIMESERIES)); // retencion
}
\end{lstlisting}

\subsection{Ejemplo de documento (registro de invitado)}
\begin{lstlisting}[style=json,caption={Documento de la colección registrations.json}]
{
  "id": "mudcvuhsnl34",
  "nombre": "Maria",
  "apellido": "Perez",
  "telefono": "099 123 4567",
  "apSsid": "Cafe-EILOR",
  "source": "esp32",
  "ts": 1790121584629
}
\end{lstlisting}

\section{Ruta de migración a base de datos relacional}
\label{sec:migracion}
Para escenarios multi\nobreakdash-dispositivo o de retención prolongada se propone migrar a un motor relacional (SQLite en despliegues embebidos o PostgreSQL en la nube). El Listado~\ref{lst:sql} presenta el esquema relacional equivalente, que preserva las entidades del modelo ER y agrega integridad referencial e índices para consultas eficientes.

\begin{lstlisting}[language=SQL,caption={Esquema relacional propuesto},label={lst:sql},basicstyle=\ttfamily\footnotesize,keywordstyle=\color{codekw}\bfseries]
CREATE TABLE dispositivo (
  mac         VARCHAR(32) PRIMARY KEY,
  tipo        VARCHAR(8)  NOT NULL,   -- 'wifi' | 'ble'
  ssid        VARCHAR(64),
  rssi        INTEGER,
  channel     INTEGER,
  vendor      VARCHAR(64),
  random_mac  BOOLEAN DEFAULT FALSE,
  evil_twin   BOOLEAN DEFAULT FALSE,
  last_seen   BIGINT
);

CREATE TABLE registro (
  id        VARCHAR(24) PRIMARY KEY,
  nombre    VARCHAR(60) NOT NULL,
  apellido  VARCHAR(60),
  telefono  VARCHAR(30),
  ap_ssid   VARCHAR(40),
  fuente    VARCHAR(16),
  ts        BIGINT NOT NULL
);
CREATE INDEX idx_registro_ts ON registro (ts);

CREATE TABLE muestra_temporal (
  t        BIGINT PRIMARY KEY,
  wifi     INTEGER, ble INTEGER,
  activos  INTEGER, paquetes BIGINT, tasa INTEGER
);
\end{lstlisting}

Esta migración es transparente para el resto del sistema gracias al desacoplamiento que ofrece \code{lib/store}: bastaría reimplementar sus funciones (\code{loadDevices}, \code{saveDevices}, \code{loadTimeseries}, \code{saveTimeseries}, \code{loadRegistrations}, \code{saveRegistrations}) sobre el nuevo motor, sin alterar \code{server.js}.
